% -------------------------------------------------------------------------------------
\section{Glauber 算符及其与光锥序证明的对照}\label{sec:scet-glauber}

标准的 SCET$_{\rm I}$、SCET$_{\rm II}$ 拉氏量里没有 Glauber 模式；Bauer--Lange--Ovanesyan 找到了必须引入 Glauber 胶子才能自洽的旁观者相互作用\cite{Bauer:2010cc}。在图形方法中 Glauber 区是 Drell--Yan 因子化最困难的一步（\cref{prop:method-glauber-deformation}，[A]；附录~O）。Glauber 交换不传播（$q^2\simeq-\bm q_{\mathrm T}^2$ 远离质壳），在有效理论中表现为横向的瞬时势；本节先陈述 Rothstein--Stewart 的构造及其地位，再给出它与附录~O 的光锥序证明之间的逐项对照。

\begin{proposition}[Glauber 算符 \tagA]\label{prop:scet-glauber-operators}
（Rothstein--Stewart\cite{RothsteinStewart2016}。）(a) 存在耦合 $n$、$\bar n$ 与软扇区的领头幂次 Glauber 算符，SCET$_{\rm II}$ 的 Glauber 拉氏量为\cite[式~(5.21)]{RothsteinStewart2016}
\begin{equation}\label{eq:scet-LG}
\mathcal L_G^{\rm II(0)}=\ee^{-\ii x\cdot\mathcal P}\sum_{n,\bar n}\sum_{i,j=q,g}\mathcal O_n^{iB}\frac{1}{\mathcal P_\perp^2}\mathcal O_s^{BC}\frac{1}{\mathcal P_\perp^2}\mathcal O_{\bar n}^{jC}
+\ee^{-\ii x\cdot\mathcal P}\sum_n\sum_{i,j=q,g}\mathcal O_n^{iB}\frac{1}{\mathcal P_\perp^2}\mathcal O_s^{j_nB},
\end{equation}
% VF-FRO-6 (re-read in the RS source: eq. (5.21))
其中 $\mathcal O_n^{qB}=\bar\chi_nT^B\frac{\bar{\sla n}}{2}\chi_n$，$\mathcal O_n^{gB}=\frac\ii2f^{BCD}\mathcal B^C_{n\perp\mu}\frac{\bar n}{2}\cdot(\mathcal P+\mathcal P^\dagger)\mathcal B^{D\mu}_{n\perp}$，$\mathcal O_{\bar n}$ 同理（$n\leftrightarrow\bar n$），$\mathcal P$ 为标签动量算符；不含软胶子时 $\mathcal O_s^{BC}=8\pi\as\,\delta^{BC}\mathcal P_\perp^2$，而对横向动量转移 $q_\perp$ 有 $\mathcal P_\perp^2\to q_\perp^2=-\bm q_{\mathrm T}^2$，故第一项在 $n$ 与 $\bar n$ 共线夸克之间产生 $\propto1/\bm q_{\mathrm T}^2$ 的势，它与 $n\cdot q$、$\bar n\cdot q$ 无关，即在 $x^\pm$ 上是局域的（瞬时的）；BPS 场重定义之前 $\mathcal L_G^{\rm I(0)}=\mathcal L_G^{\rm II(0)}$。
(b) 在硬散射中，Glauber 交换可以引起使因子化失效的修正；对硬散射算符，软与共线圈图的 Glauber 零箱减除保证结果对软与共线 Wilson 线的方向不敏感\cite[摘要]{RothsteinStewart2016}。
(c) 对旁观者—旁观者交换，迭代的 Glauber 交换给出一个相位，当横向动量转移 $\Delta p_\perp$ 的积分不受限制时，它在截面中相消。
(d) 对活跃—活跃交换（SCET$_{\rm II}$），软 Wilson 线取与被探测的共线粒子的物理方向一致的方向时，Glauber 交换图等于软图的 Glauber 零箱，二者可以一并吸收进软 Wilson 线\cite[第~10、10.1~节]{RothsteinStewart2016}。
（本命题只涉及部分子散射。）
\uses{def:scet-modes, prop:scet-scet2}
\end{proposition}

\begin{proofsketch}
% errtype: E-S
(1) \emph{树图匹配。}$t$ 道单胶子交换的前向散射在领头幂次中，$\bar u_n\gamma^\mu u_n=n^\mu\,\bar u_n\frac{\bar{\sla n}}{2}u_n+\order{\lambda}$（$\bar n$ 共线夸克同理），而 Glauber 动量使 $q^2=-\bm q_{\mathrm T}^2[1+\order{\lambda^2}]$，故
\[
\ii\mathcal M=-\ii\,\frac{8\pi\as}{\bm q_{\mathrm T}^2}\Big[\bar u_n\frac{\bar{\sla n}}{2}T^Bu_n\Big]\Big[\bar u_{\bar n}\frac{\sla n}{2}T^Bu_{\bar n}\Big]+\order{\lambda},
\]
取不含软胶子的 $\mathcal O_s^{BC}=8\pi\as\delta^{BC}\mathcal P_\perp^2$，这正是\cref{eq:scet-LG}第一项的矩阵元（数值上相对偏差在 $\lambda=10^{-1},10^{-2},10^{-3}$ 时为 $3.2\times10^{-3}$、$2.8\times10^{-5}$、$2.8\times10^{-7}$）；含软胶子的部分由带一个与两个软胶子的匹配确定，其单胶子 Feynman 规则是 Lipatov 顶角\cite[第~5.1、6~节]{RothsteinStewart2016}。% w18_s5 (D1, D2)
(2) \emph{幂次计数。}Glauber 算符是领头幂次的，且其形式没有圈修正\cite[第~5.3~节]{RothsteinStewart2016}（E-S 型陈述）。
(3) \emph{Glauber 圈。}两条共线线交换两个 Glauber 势（圈动量 $k$）时，领头幂次的两个分母分别只含 $n\cdot k=k^0-k^z$ 与 $\bar n\cdot k=k^0+k^z$；每个 Glauber 势带一个快度正规子因子 $|2k^z|^{-\eta}\nu^\eta$，它只依赖 $k^z$，故可先作 $k^0$ 积分。交叉箱图中两个 $k^0$ 极点在实轴同侧，积分精确为零；箱图中两个极点分居两侧，余下的 $k^z$ 积分的主值部分 $\propto\tan(\pi\eta)$ 在 $\eta\to0$ 时趋于零，只剩来自 $\delta$ 函数的纯虚部分（在 Rothstein--Stewart 的归一化下为 $-\ii\pi$），这正是 Glauber 势指数化为相位所需的因子\cite[第~5.2.1~节，式~(5.30)、(5.31)]{RothsteinStewart2016}。% w18_s5 (E1-E3), w18_s6 (B0-B2)
(4) \emph{硬匹配。}在单圈与两圈显式计算中，活跃线之间的 Glauber 交换给出全部低能的 $\ii\pi$ 项，并与软图的 Glauber 零箱对应，由此得 (b)、(d)\cite[第~10~节]{RothsteinStewart2016}。
(5) \emph{旁观者。}Glauber 梯形图之和是相位 $\ee^{\ii\phi(b_\perp)}$，它在 $\Delta p_\perp$ 的积分不受限制时相消，得 (c)\cite[式~(9.1)、(9.18)--(9.19)、(11.12)--(11.13)]{RothsteinStewart2016}；同一机制见\cite[式~(14)]{AybatSterman2009}。完整论证见\cite[第~5、9--11~节]{RothsteinStewart2016}。
\end{proofsketch}

\paragraph{地位与缺口}
算符的构造（(a)）是由匹配与幂次计数得到的全阶论证；(b)--(d) 是在单圈、两圈与梯形图求和中显式得到的，作者认为这一对应“very plausible”（非常可信）地推广到所有阶，但没有给出证明\cite[第~10.2~节]{RothsteinStewart2016}。作者明确指出旁观者相位的相消“taken alone it is far from a proof of factorization, even in the abelian case”（单独来看远不是因子化的证明，即使在阿贝尔情形）\cite[第~11.1~节]{RothsteinStewart2016}，而 SCET 中完整的因子化证明需要证明 Glauber 拉氏量对硬散射没有贡献，并称这一证明将在别处给出\cite[第~11~节]{RothsteinStewart2016}——据我们截至 2026 年 10 月的 INSPIRE 文献检索，尚未见到。因此最弱的一步是“$\mathcal L_G$ 的效应在单举截面中相消”，它对应附录~O 的\cref{asm:appO-mixing}(c) 与\cref{asm:appO-deformation}；这是相消问题而不是幂次估计（作为幂次计数，(a) 只给出 E-S 型陈述）。对 $e^+e^-$ 的 $n$–$\bar n$ 产生流，作者在单圈与两圈显式表明：在 SCET$_{\rm II}$ 中活跃线之间的 Glauber 贡献等于软图的 Glauber 零箱，在 SCET$_{\rm I}$ 中 Glauber 图不贡献、$\ii\pi$ 项由超软图携带，并预期在任意阶与不含 Glauber 交换的理论结果相同\cite[第~10.1、10.3~节]{RothsteinStewart2016}——这只是单圈、两圈的佐证。在强子对撞中，Gaunt 论证了对强子 $E_T$ 而言 CSS 的相消论证在两个旁观者 Glauber 胶子处失效\cite{Gaunt2014}，Rothstein--Stewart 也认为当测量约束 $\Delta p_\perp$ 时旁观者交换从 $\order{\as^4}$ 起破坏因子化\cite[第~11.1~节]{RothsteinStewart2016}；Schwartz--Yan--Zhu 对纯 Glauber 梯形图与任意红外安全的单标度观测量给出了全阶相消的论证，否定了这些论断中纯 Glauber 的部分，但不适用于多重微分的观测量\cite{Schwartz:2018obd}（第~18、24~章，\cref{prop:front-glauber-bfkl,prop:front-syz}）。升级途径：在有效理论中证明 $\mathcal L_G$ 对单举观测量的贡献相消，即\cref{thm:appO-main}（[R, $R_{\mathrm{cond}}$]）去掉模型限制后的有效理论版本。

\begin{remark}[Glauber SCET 与光锥序证明的对照]\label{rem:scet-glauber-dictionary}
Rothstein--Stewart 的 Glauber 有效理论与附录~O 的光锥序（LCOPT）证明处理的是同一个动量区域（精确版本见\cref{rem:appO-rs-dictionary}）：
\begin{center}\small
\renewcommand{\arraystretch}{1.1}
\begin{tabular}{@{}p{0.4\textwidth}p{0.53\textwidth}@{}}
\toprule
Glauber SCET\cite{RothsteinStewart2016} & 附录~O 的光锥序证明 \\
\midrule
瞬时势 $\mathcal O_n\frac{1}{\mathcal P_\perp^2}\mathcal O_s\frac{1}{\mathcal P_\perp^2}\mathcal O_{\bar n}$ & 等光锥时刻插入与 Glauber 算符 $V_G$（\cref{lem:appO-equal-time,lem:appO-glauber-operator}） \\
Glauber 的“塌缩”到硬顶点 & 与 $V_H$ 重合的接触项（\cref{rem:appO-contact}） \\
旁观者相位在 $\Delta p_\perp$ 积分后相消 & 最大时刻配对、族求和、窗口与整函数载子（\cref{def:appO-windows,lem:appO-carrier}） \\
Glauber 零箱；Glauber 与软模式的混合 & 减除法中的零箱（\cref{def:scet-zero-bin}，附录~N）；\cref{asm:appO-mixing} \\
快度正规子 $|2k^z|^{-\eta}$：交叉箱图为零 & 共同时序，$p_{\rm opp}=0$（\cref{rem:dy-common-ordering}；\cref{lem:appO-glauber-operator}(c)） \\
\bottomrule
\end{tabular}
\end{center}
最后一行的机制已在\cref{prop:scet-glauber-operators}证明概要的第~(3)~步中给出：正规子只依赖 $k^z$，先作 $k^0$ 积分，交叉箱图精确为零，箱图只给出相位。在附录~O 的语言中，这正是全部 $s$ 道权重落在非交叉次序上的“共同时序”，它是精确图给出的物理结果，独立的对称截断给出的 $p_{\rm opp}=\frac12$ 只是截断假象（\cref{rem:dy-common-ordering}）。第四行的对应也有文献依据：Rothstein--Stewart 指出，硬散射中活跃—活跃的 Glauber 贡献可以吸收进软 Wilson 线，这与 CSS 的围道变形图像一致——软与 Glauber 的合并圈积分被变形离开 Glauber 区\cite[第~10.1~节]{RothsteinStewart2016}。
\end{remark}
